\subsection{Recognition and memory}
\label{sec:results-decon}
% T2b: the treated (decontaminated) text beside the original, means over cutoffs. Read only under
% \nwdecontrue, from sections/06b-decon.tex. Keys: _build/fill_decon.py.
\begin{table*}[tp]
\centering\small
\setlength{\tabcolsep}{2.6pt}
\begin{tabular}{@{}llcccccccc@{}}
\toprule
 &  & \multicolumn{3}{c}{\nwhead{Recall}} & \multicolumn{2}{c}{\nwhead{Precision}} & \multicolumn{2}{c}{\nwhead{Citation sufficiency}} & \nwhead{Later facts} \\
\cmidrule(lr){3-5}\cmidrule(lr){6-7}\cmidrule(lr){8-9}
\multicolumn{2}{@{}l}{\nwhead{System}} & \nwhead{treated} & \nwhead{original} & \nwhead{difference} & \nwhead{treated} & \nwhead{original} & \nwhead{treated} & \nwhead{original} & \nwhead{disclosed, treated} \\
\midrule
B1 & Long context & \cellcolor{nwhighcyan}\best{\res{M1-annedecon-avg-B1}}\sdev{\res{M1-annedecon-sd-B1}} & \respt{M1-anne-avg-B1} & \res{M1-annedecon-diff-B1} & \res{M3-annedecon-avg-B1}\sdev{\res{M3-annedecon-sd-B1}} & \respt{M3-anne-avg-B1} & \res{M8-annedecon-avg-B1}\sdev{\res{M8-annedecon-sd-B1}} & \respt{M8-anne-avg-B1} & \res{M4-annedecon-total-B1} \\
B2 & Prefix retrieval & \res{M1-annedecon-avg-B2}\sdev{\res{M1-annedecon-sd-B2}} & \respt{M1-anne-avg-B2} & \res{M1-annedecon-diff-B2} & \res{M3-annedecon-avg-B2}\sdev{\res{M3-annedecon-sd-B2}} & \respt{M3-anne-avg-B2} & \res{M8-annedecon-avg-B2}\sdev{\res{M8-annedecon-sd-B2}} & \respt{M8-anne-avg-B2} & \res{M4-annedecon-total-B2} \\
B3 & LightRAG & \res{M1-annedecon-avg-B3}\sdev{\res{M1-annedecon-sd-B3}} & \respt{M1-anne-avg-B3} & \res{M1-annedecon-diff-B3} & \cellcolor{nwhighcyan}\best{\res{M3-annedecon-avg-B3}}\sdev{\res{M3-annedecon-sd-B3}} & \respt{M3-anne-avg-B3} & \res{M8-annedecon-avg-B3}\sdev{\res{M8-annedecon-sd-B3}} & \respt{M8-anne-avg-B3} & \res{M4-annedecon-total-B3} \\
B4 & Graph baseline & \res{M1-annedecon-avg-B4}\sdev{\res{M1-annedecon-sd-B4}} & \respt{M1-anne-avg-B4} & \res{M1-annedecon-diff-B4} & \res{M3-annedecon-avg-B4}\sdev{\res{M3-annedecon-sd-B4}} & \respt{M3-anne-avg-B4} & \cellcolor{nwhighcyan}\best{\res{M8-annedecon-avg-B4}}\sdev{\res{M8-annedecon-sd-B4}} & \respt{M8-anne-avg-B4} & \cellcolor{nwhighpink}\best{\res{M4-annedecon-total-B4}} \\
B5 & \sys & \res{M1-annedecon-avg-B5}\sdev{\res{M1-annedecon-sd-B5}} & \respt{M1-anne-avg-B5} & \res{M1-annedecon-diff-B5} & \res{M3-annedecon-avg-B5}\sdev{\res{M3-annedecon-sd-B5}} & \respt{M3-anne-avg-B5} & \res{M8-annedecon-avg-B5}\sdev{\res{M8-annedecon-sd-B5}} & \respt{M8-anne-avg-B5} & \cellcolor{nwhighpink}\best{\res{M4-annedecon-total-B5}} \\
\midrule
X1 & Closed book & \res{M1-annedecon-avg-X1}\sdev{\res{M1-annedecon-sd-X1}} & \respt{M1-anne-avg-X1} & \cellcolor{nwhighpink}\best{\res{M1-annedecon-diff-X1}} & \res{M3-annedecon-avg-X1}\sdev{\res{M3-annedecon-sd-X1}} & \respt{M3-anne-avg-X1} & -- & -- & \res{M4-annedecon-total-X1} \\
X2 & Full text & \res{M1-annedecon-avg-X2}\sdev{\res{M1-annedecon-sd-X2}} & \respt{M1-anne-avg-X2} & \res{M1-annedecon-diff-X2} & \res{M3-annedecon-avg-X2}\sdev{\res{M3-annedecon-sd-X2}} & \respt{M3-anne-avg-X2} & \res{M8-annedecon-avg-X2}\sdev{\res{M8-annedecon-sd-X2}} & \respt{M8-anne-avg-X2} & \cellcolor{nwhighpink}\best{\res{M4-annedecon-total-X2}} \\
\bottomrule
\end{tabular}
\caption{Original and treated text (names replaced, paragraphs paraphrased), with unchanged
code, prompts and gold facts. Rates are unweighted means over $t=1..5$; treated values
include sample standard deviations. Differences: treated minus original, before rounding.
Disclosure counts cover $t=1..4$.
Claude finalized the AI labels in both conditions. Bold highlights:
best treated B1--B5 rates (cyan), largest recall loss and nonzero disclosures (pink).}
\label{tab:decon}
\end{table*}

All seven systems were rebuilt on renamed, paraphrased text, retaining paragraph identifiers
to align gold facts and citations (Appendix~\ref{app:contamination}). In 200 aligned probes,
title identification falls from \res{MEM-title-orig} to \res{MEM-title-decon}; a names-only
control gives \res{MEM-title-remap} (Table~\ref{tab:recognition}). Residual recognition
and imperfect fidelity limit interpretation.

Closed-book recall falls from \respt{M1-anne-avg-X1} to
\respt{M1-annedecon-avg-X1} (rounded). One gold fact--cutoff match remains: a death in an
invented story (\res{M1-annedecon-count-X1} assessments). Of 200 sampled closed-book assertions,
\res{M3-annedecon-unsup-X1} are unsupported. Text-based systems retain much of their recall,
with a maximum mean loss of \res{M1-annedecon-maxloss} (Table~\ref{tab:decon}).
The contrast is consistent with using supplied evidence, but does not isolate each system's
memory contribution. Long-context citation sufficiency declines from
\respt{M8-annedecon-t1-B1} at $t=1$ to \respt{M8-annedecon-t5-B1} at $t=5$, often because
citations retain a paragraph number but omit its volume.

B1, B2, B3 and X1 disclose no assessed later fact on the treated text. \sys states the
leads' marriage at $t=2$ and $t=3$ (\res{M4-annedecon-total-B5}): a neighbor's prediction
was extracted as a presumed marriage and rendered as a spouse relation. B4, using the
same graph, implies it (\res{M4-annedecon-total-B4}). The failure concerns attribution
and modality in eligible evidence, which a volume filter cannot repair. X2 discloses
\res{M4-annedecon-total-X2} later facts, versus \res{M4-anne-total-X2} on the original;
\res{M4-annedecon-t1-X2} occur at $t=1$, where the original run had none. The full-text
instruction remains unreliable after recognition is reduced.
